% !TeX root = ../main.tex

\begin{abstract}

隨著網路食譜平台的蓬勃發展，使用者面臨嚴重的資訊過載問題，亟需有效的推薦系統協助探索潛在偏好的食譜內容。本研究以知識圖譜注意力網路 (Knowledge Graph Attention Network, KGAT) 為核心，在 Food.com 食譜推薦資料集上進行系統性的效能評估與可解釋性研究。

本研究的主要貢獻包含三個面向。首先，透過消融實驗，顯示圖神經網路的訊息傳播深度是影響推薦效能的重要因素：三層架構 ($L{=}3$) 的 HR@20 達 0.8775，較單層架構 ($L{=}1$) 提升逾 29\%，較最強的基準方法 LightGCN 更提升達 31\%，全面超越 BPR-MF、LightGCN 及 NFM 等基準方法。其次，發現在淺層架構 ($L{=}1$) 下，引入知識圖譜反而因雜訊稀釋協同過濾訊號而導致效能下降，此發現與近期文獻~\cite{zhang2024kg4receval}所觀察到之知識圖譜貢獻非必然正面的結論相呼應，為知識圖譜整合策略提供了重要啟示。最後，本研究以 500 位使用者為基礎，對三種不同深度的 KGAT 模型（$L{=}1$、$L{=}2$、$L{=}3$）進行跨模型 Fidelity 量化評估。在最佳效能模型 KGAT-3L 中，278 位使用者成功萃取解釋路徑（覆蓋率 55.6\%），其中 94.2\% 的路徑 Fidelity$^+$ 大於零，支持注意力權重標記了對預測有貢獻的關鍵邊。進一步的路徑類型分析發現，直接連接路徑的忠實度最高（Avg $F^+=0.0494$，$F^- \approx 0$），而 KG 語意路徑的充分性忠實度持續偏低（$F^-=0.2449$），暗示深層模型的決策可能已將部分知識圖譜結構資訊壓縮進分散式嵌入中，使得顯式圖譜路徑無法完整涵蓋模型的決策依據。

\end{abstract}

\begin{abstract*}

With the rapid growth of online recipe platforms, users face severe information overload, necessitating effective recommender systems to help discover recipes matching their preferences. This study systematically evaluates the performance and explainability of the Knowledge Graph Attention Network (KGAT) on the Food.com recipe recommendation dataset.

Our contributions are threefold. First, through ablation experiments, we show that the message propagation depth of graph neural networks is an important factor for recommendation performance: the three-layer architecture ($L{=}3$) achieves an HR@20 of 0.8775, representing a 29\% improvement over the single-layer architecture ($L{=}1$) and a 31\% improvement over the strongest baseline LightGCN, substantially outperforming baseline methods including BPR-MF, LightGCN, and NFM. Second, we discover that incorporating the knowledge graph in the shallow architecture ($L{=}1$) paradoxically degrades performance, as noise dilutes the collaborative filtering signals---a finding consistent with recent work by Zhang et al.~\cite{zhang2024kg4receval} questioning the unconditional benefit of knowledge graphs in recommender systems---providing important insights for knowledge graph integration strategies. Third, using a cohort of 500 users, we conduct a cross-model Fidelity evaluation across three KGAT depths ($L{=}1$, $L{=}2$, $L{=}3$). In the best-performing KGAT-3L model, 278 users yield explanation paths (coverage 55.6\%), of which 94.2\% exhibit positive Fidelity$^+$, supporting that attention weights identify edges contributing to predictions. Path-type analysis further reveals that direct connection paths achieve the highest fidelity (Avg $F^+=0.0494$, $F^- \approx 0$), while KG semantic paths exhibit persistently poor sufficiency fidelity ($F^-=0.2449$), suggesting that deep models may partly compress knowledge graph structural information into distributed representations where explicit graph paths cannot fully capture the model's decision evidence.

\end{abstract*}